\documentclass[10pt,a4paper]{amsart}
\usepackage[margin=19mm]{geometry}
\usepackage{amsmath,amssymb,mathtools}
\usepackage[scr=rsfs,cal=euler]{mathalfa}
\usepackage{xfrac}
\usepackage[unicode,hidelinks,bookmarks=false]{hyperref}
\newcommand{\ie}{i.e.\ }
\newcommand{\cf}{cf.\ }
\newcommand{\resp}{respectively\ }
\newcommand{\Subsection}{\subsection}
\providecommand{\nobreakdash}{\nobreak\hbox{-}}
\setlength{\emergencystretch}{2em}
\multlinegap0pt
\usepackage{mathtools}
\usepackage{tikz}
\usepackage{xcolor}
\usepackage{amssymb}
\newtheorem{theorem}{Theorem}
\newtheorem{lemma}{Lemma}
\newcommand{\calF}{{\mathcal{F}}}
\newcommand{\calO}{{\mathcal{O}}}
\newcommand{\reg}{\operatorname{reg}}
\title{Multigraded Regularity of the Complete Flag Variety}
\author{Caitlin M. Davis}
\date{Source: arXiv:2606.21058v1 (CC BY 4.0)}
\begin{document}
\maketitle

\noindent\emph{Source and licence.} Multigraded Regularity of the Complete Flag Variety. Version arXiv:2606.21058v1 (CC BY 4.0), distributed by arXiv under the Creative Commons Attribution 4.0 International licence, http://creativecommons.org/licenses/by/4.0/, as recorded in the arXiv licence metadata for this exact version and shown on the abstract page https://arxiv.org/abs/2606.21058v1. Full licence notice: \texttt{licenses/arxiv-2606.21058-CC-BY-4.0.txt}.\par\smallskip

\noindent\textit{Editorial scope.} Counted unit: the article's lemma lemma:reverse (source0.tex lines 522--525). The article's complete original proof of it is delivered with it (source0.tex lines 527--564, one \texttt{proof} environment of 38 lines). No mathematics is written, repaired or re-derived: the statement and the proof are the article's own lines, in the article's own notation, and no retained line is substituted or re-flowed.  Retained uncounted as the source-local context the proof needs: lines 350--356.  Nothing was omitted from any retained span, so the package carries no omission record.  No retained line contains a citation, so the wrapper carries no bibliography and no bibliography entry.  No retained line contains a figure environment, an image-inclusion command or drawing code, so no image is delivered and the package declares no asset dependency.  Editorial formatting only, in addition: an AMS \texttt{amsart} preamble that restates the article's own declarations and macros used by the retained lines, namely the environments \texttt{theorem}, \texttt{lemma} on separate counters, together with the standard packages those lines need.  The retained environments print in this document as Theorem 1, Lemma 1 in order of appearance, whereas the article prints the same items with the numbering of its own full document.  The complete native source of the article as distributed in the arXiv e-print for this exact version is bundled unchanged as \texttt{sources/203254/source0.tex}.
\begin{theorem}[Bott's theorem]\label{thm:bottAlternative}
    Consider $\mathbf b \coloneqq \mathbf a - (0,1,2,\ldots,n-1)$.  If two entries of $\mathbf b$ are equal, then $H^i(\calF,\calO_{\calF}(\mathbf a)) = 0 $ for all $i$.  Otherwise, let $N$ be the number of inversions in $\mathbf b$.  That is, $N$ is the number of pairs $i<j$ such that $b_i<b_j$.  Then
    \begin{align*}
    &H^N(\calF,\calO_{\calF}(\mathbf a)) \neq 0, \text{ and } \\
    &H^i(\calF,\calO_{\calF}(\mathbf a)) = 0 \text{ for } i\neq N.
    \end{align*}
\end{theorem}

\begin{lemma}
\label{lemma:reverse}
    Let $\mathbf d \in \mathbb N^{n-1}$ and let $\overline{\mathbf d} = (d_{n-1},\ldots,d_1)$ be given by reversing the entries of $\mathbf d$.  Then $\mathbf d$ is in $\reg(\iota_*\calO_{\calF})$ if and only if $\overline{\mathbf d}$ is in $\reg(\iota_*\calO_{\calF})$.
\end{lemma}

\begin{proof}
    We only need to argue one direction, so suppose $\mathbf d \notin \reg(\iota_*\calO_{\calF})$. Then there is some $i \leq \dim \calF$ and some $\mathbf c \in \mathbb N^{n-1}$ such that $|\mathbf c| = i$ and 
    \begin{equation}\label{eqn:nonzeroCohom}
        H^i(\mathbb P,\iota_*\calO_{\calF}\otimes \calO_{\mathbb P}(\mathbf d-\mathbf c)) \neq 0
        %H^i(\calF, \iota^*\calO_{\mathbb P}(\mathbf{a})\otimes \iota^*\calO_{\mathbb P}(-\mathbf{f})) \neq 0.
    \end{equation}
    
    For ease of notation, write $\mathbf S = \varphi(\mathbf d - \mathbf c)$.  Then (\ref{eqn:nonzeroCohom}) becomes:
    \begin{equation}\label{eqn:nonzeroCohom2}
        H^i(\calF,\calO_{\calF}(\mathbf S)) \neq 0.
        %H^i(\calF,\calO_{\calF}(S_1,\ldots,S_n)) \neq 0.
    \end{equation}
    By Theorem \ref{thm:bottAlternative}, this means $\mathbf S - (0,1,\ldots,n-1)$ contains exactly $i$ inversions.  
    
    We will argue that
    \begin{equation}\label{eqn:NTSCohom}
        H^i(\mathbb P,\iota_*\calO_{\calF}\otimes \calO_{\mathbb P}(\overline{\mathbf d}-\overline{\mathbf c})) \neq 0
        %H^i(\calF, \iota^*\calO_{\mathbb P}(\mathbf{\overline a})\otimes \iota^*\calO_{\mathbb P}(-\mathbf{\overline f})) \neq 0,
    \end{equation}
    so that $\mathbf{\overline d} \notin \reg(\iota_*\calO_{\calF})$.  We can rewrite (\ref{eqn:NTSCohom}) as:
    \[
    H^i(\calF,\calO_{\calF}(\varphi(\overline{\mathbf d}-\overline{\mathbf c}))) \neq 0.
    \]
    We compute:
    \begin{align*}
        \varphi(\overline{\mathbf d}-\overline{\mathbf c}) &= ((d_1-c_1))+\ldots+(d_{n-1}-c_{n-1}),\ldots, (d_1-c_1)+(d_2-c_2),(d_1-c_1),0)\\
        &= (S_1-S_n,S_1-S_{n-1},\ldots,S_1-S_3,S_1-S_2,0).
    \end{align*}
    Concretely, we need to argue that the following contains $i$ inversions:
    \[
    (S_1-S_n,S_1-S_{n-1},\ldots,S_1-S_3,S_1-S_2,0) - (0,1,\ldots,n-1).
    \]
    The number of inversions is unaffected by adding a multiple of $(1,1,\ldots,1)$, so we can add $(-S_1,\ldots,-S_1)$ to obtain:
    \[
    (-S_n,-S_{n-1},\ldots,-S_1) - (0,1,\ldots,n-1) = (-S_n - 0, -S_{n-1}-1,\ldots, -S_1-(n-1)).
    \]
    The assumption (\ref{eqn:nonzeroCohom2}) tells us concretely that there are $i$ pairs $j<k$ for which $S_j-(j-1)<S_k-(k-1)$.  This inequality holds if and only if $-S_j+j > -S_k + k$, or equivalently $-S_j-(n-j) > -S_k-(n-k)$.  Therefore, $(-S_n - 0, -S_{n-1}-1,\ldots, -S_1-(n-1))$ contains $i$ inversions, proving (\ref{eqn:NTSCohom}).
\end{proof}

\end{document}
